Lab Notes (04.22.2026): 
\begin{itemize}
    \item Connect SPAD A, SPAD B, and Laser Sync to Swabian channel 1,2 and 3 respectively. 
    \item Set Pico-Quant LASER at $70\%$ power. 
    \item LASER frequency is configured at $1MHz$. Select internal oscillator\textit{Source: Int2}, and Rep Frequency at \textit{1}.
    \item Measure Power at $20.40 \mu W$
    \item Align the LASER to the SPADs active areas. 
    \item Place ND0.6 Filter.
    \item Isolate the system from external light sources.
    \item Calibrate. Set baseline for SPAD dark counts and Laser clock. Laser head remains off.  
    \begin{itemize}
        \item Use \texttt{tt.Countrate} to extract average counts-per-second (cps) over $5*10^{11}$ picoseconds. 
        \item Averaging the previus averages 55 times, we get:
        \begin{itemize}
            \item Mean SPAD A: 7.931 cps
            \item Mean SPAD B: 9353.948 cps
            \item Mean LASER Sync: 999997.776 cps
        \end{itemize} 
    \end{itemize}
    \item Turn Laser On and calibrate cps. 
    \begin{itemize}
        \item Use \texttt{tt.Countrate} to extract average counts-per-second (cps) over $5*10^11$ picoseconds. 
        \item Averaging the previus averages 94 times, we get:
        \begin{itemize}
            \item Mean SPAD A: 170332.128 cps
            \item Mean SPAD B: 123889.021 cps
            \item Mean LASER: 999997.468cps
        \end{itemize}
    \end{itemize}
    \item Then, we proceed with the calibration for the optical path differences. We use \texttt{tt.Correlation} function, with $10$ps resolution to find the time differences between two channels and obtained the following results. 
    \begin{itemize}
        \item Sync $\rightarrow$ SPAD A Delays:
        \item[] $ [-10200, -10210, -10200 ]$ ps
        \item Sync $\rightarrow$ SPAD B Delays:
        \item[] $ [-15220, -15220, -15230] $ ps
        \item SPAD A $\rightarrow$ SPAD B Delay: 
        \item[] [$-4910, -4930, -4920] $ ps
        \item See $optical\_path\_calibration.png$
    \end{itemize}
    \item The values obtained are taken into account for the first coincidence measurement of $g^{(2)}(0)$ described in the experiment procedure in subsection \ref{subsec:CorrVal}. 
    \item We observe $g^{(2)}(0)$ value is higly dependent on the lenght of coincidence window. For 400 ps $g^{(2)}(0)=0.0001$ while for 100,000 ps we get $g^{(2)}(0)=0.9442$
    \item \textbf{Aumenting the LASER Frequency} %%%%%%%%%%%%%%%%%%%%%%%%%%%%%%%%%%%%%%%%%%%%%%%%%%%%%%%
    \item LASER frequency is configured at $80MHz$. Select internal oscillator\textit{Source: Int1}, and Rep Frequency at \textit{1}
    \item Turn Laser On and calibrate cps. 
    \begin{itemize}
        \item Use \texttt{tt.Countrate} to extract average counts-per-second (cps) over $5*10^11$ picoseconds. 
        \item Averaging the previus averages 31 times, we get:
        \begin{itemize}
            \item Mean SPAD A: 11288471.685 cps
            \item Mean SPAD B: 7219733.294 cps
            \item Mean LASER: 79999884.439 cps
        \end{itemize}
    \end{itemize}
    \item Then, we proceed with the calibration for the optical path differences. We use \texttt{tt.Correlation} function, with $10$ps resolution to find the time differences between two channels and obtained the following results. 
    \begin{itemize}
        \item Sync $\rightarrow$ SPAD A Delays:
        \item[] $ [2445, 2440  ]$ ps
        \item Sync $\rightarrow$ SPAD B Delays:
        \item[] $ [-2580, -2590 ] $ ps
        \item SPAD A $\rightarrow$ SPAD B Delay: 
        \item[] $[-4915, -4920] $ ps
        \item See $optical\_path\_calibration.png$
    \end{itemize}
    \item The values obtained are taken into account for the first coincidence measurement of $g^{(2)}(0)$ described in the experiment procedure in subsection \ref{subsec:CorrVal}. 
    \item We observe $g^{(2)}(0)$ value is higly dependent on the lenght of coincidence window. For 400 ps $g^{(2)}(0)=0.0007$ while for 2,000 ps we get $g^{(2)}(0)=0.9836$.
\end{itemize}
